\chapter{Συζήτηση και Συμπεράσματα}
\label{chap:conclusions}

Το κεφάλαιο αυτό ερμηνεύει τα αποτελέσματα της εργασίας. Τα αντιπαραβάλλει με τους τρεις στόχους του επίσημου κειμένου ανάθεσης και καταγράφει τους περιορισμούς και τις κατευθύνσεις μελλοντικής βελτίωσης. Η εργασία παρέχει μια αναπαραγώγιμη βάση για τη μελέτη αλγορίθμων βέλτιστης επίλυσης του κύβου του Rubik. Η υλοποίηση καλύπτει την αναπαράσταση του κύβου 3×3×3, τον έλεγχο εγκυρότητας, την ακριβή ρηχή αναζήτηση και τον επιλυτή Korf/IDA* (\eng{iterative-deepening A*}) με εγγενείς πίνακες προτύπων (\eng{pattern databases}, PDB). Περιλαμβάνει επίσης τη μηχανή H48 ακριβούς αναζήτησης, τους επιλυτές Kociemba και Thistlethwaite, τον πρακτικό προσαρμογέα (\eng{adapter}) δύο φάσεων και την πλήρη εξαντλητική μελέτη του κύβου 2×2×2 (\eng{Pocket Cube}). Όλα τα αποτελέσματα αποθηκεύονται ως τεκμήρια και επαληθεύονται πριν χρησιμοποιηθούν στο κείμενο, όπως περιγράφεται στο Κεφάλαιο~\ref{chap:validation}.

\section{Συζήτηση}
\label{sec:conclusions-discussion}

Το κεντρικό ερμηνευτικό συμπέρασμα της εργασίας είναι η ουσιώδης διαφορά ανάμεσα στο «λύνω» και στο «αποδεικνύω βέλτιστο μήκος». Μια ακολουθία κινήσεων μπορεί να είναι ορθή, σύντομη και πρακτικά χρήσιμη χωρίς να είναι αποδεδειγμένα ελάχιστη· πρόκειται για διαφορετικό τύπο ισχυρισμού και όχι για μειονέκτημα της πρακτικής επίλυσης. Για τον λόγο αυτό, η ετικέτα \code{exact} (αποδεδειγμένα βέλτιστη λύση) αποδίδεται μόνο όταν μια πλήρης ή παραδεκτή αναζήτηση έχει ολοκληρωθεί: η αναζήτηση κατά πλάτος (\eng{breadth-first search}, BFS), ο αλγόριθμος IDA* και το υποσύστημα H48 δίνουν ακριβείς απαντήσεις μόνο όπου η αντίστοιχη αναζήτηση ολοκληρώνεται. Αντίθετα, οι λύσεις των Kociemba και Thistlethwaite, αν και επαληθεύονται σε ολόκληρο το σταθερό σώμα μετρήσεων, χαρακτηρίζονται \code{non_exact} (επαληθευμένη λύση χωρίς απόδειξη βελτιστότητας), επειδή δεν συνοδεύονται από απόδειξη ελαχίστου μήκους. Η ετικέτα αυτή δεν δηλώνει αποτυχία· αποτελεί έντιμη περιγραφή της παρεχόμενης εγγύησης.

Με την ίδια λογική, οι εξαντλήσεις χρονικού ορίου (\code{timeout}, εξάντληση του προϋπολογισμού αναζήτησης χωρίς λύση) και τα κάτω φράγματα (\code{lower_bound}, αποδεδειγμένο κάτω φράγμα χωρίς πλήρη λύση) είναι πειραματικά αποτελέσματα και όχι κενά προς απόκρυψη. Η ρητή καταγραφή τους προστατεύει από δύο σφάλματα ερμηνείας. Το πρώτο θα ήταν να θεωρηθεί ότι, επειδή ένας άλλος επιλυτής βρήκε λύση, η ακριβής αναζήτηση θα έδινε αναγκαστικά την ίδια απόσταση· το δεύτερο, ότι η εξάντληση χρονικού ορίου σημαίνει μη επιλυσιμότητα. Καμία από τις δύο ερμηνείες δεν ισχύει: η εξάντληση χρονικού ορίου δηλώνει μόνο ότι η συγκεκριμένη διαμόρφωση δεν ολοκληρώθηκε εντός του ορίου της. Στο κύριο σώμα μετρήσεων, εξαντλήσεις χρονικού ορίου παραμένουν μόνο στον επιλυτή Korf/IDA* καθαρής Python, για τα βαθύτερα τυχαία ανακατέματα (\eng{scramble})· οι εγγενείς υλοποιήσεις με πίνακες προτύπων και το υποσύστημα H48 αποδεικνύουν προοδευτικά βαθύτερες περιπτώσεις ως ακριβείς. Τα επιμέρους μεγέθη και η κλιμάκωση του χρόνου με το βάθος παρουσιάζονται στο Κεφάλαιο~\ref{chap:experiments} και στο Σχήμα~\ref{fig:runtime-depth}.

Οι πίνακες προβολών, η γωνιακή PDB και οι εξακμικοί πίνακες προτύπων συνεισφέρουν σε τρία επίπεδα. Πρώτον, δίνουν πραγματικό, καθοδηγούμενο από πίνακες κάτω φράγμα στον επιλυτή Korf/IDA*. Δεύτερον, τεκμηριώνουν τους στόχους των φάσεων Kociemba και Thistlethwaite, επειδή οι στόχοι αυτοί συνδέονται με συγκεκριμένες συντεταγμένες. Τρίτον, παράγουν μετρήσιμα τεκμήρια με μεγέθη πεδίων, πλήθη εγγραφών και αθροίσματα ελέγχου, ώστε η έννοια της ευρετικής να μη μένει αφηρημένη. Η ισχύς τους, ωστόσο, είναι περιορισμένης εμβέλειας: οι προβολές CO/EO/UD-slice αγνοούν μεγάλο μέρος της κατάστασης, η γωνιακή PDB αγνοεί τις ακμές, κάθε εξακμικός πίνακας βλέπει μόνο έξι ακμές και ο κύριος συνδυασμός τους γίνεται με μέγιστο. Τα πειράματα επιμερισμένου κόστους (\eng{cost-partitioned}) δείχνουν ότι η υλοποίηση υποστηρίζει ασφαλές προσθετικό άθροισμα όταν τα κόστη των τελεστών διαμερίζονται ρητά. Η συγκεκριμένη κατανομή URF/DLB, όμως, είχε μέγιστη προβεβλημένη απόσταση 3 και δεν βελτίωσε το ντετερμινιστικό δείγμα πέρα από το μέγιστο των οκτώ εξακμικών πινάκων πλήρους κόστους. Γι' αυτό ο επιλυτής Korf/PDB έχει περιορισμένη εμβέλεια: η υποδομή αρκεί για ουσιαστική εγγενή ακριβή αναζήτηση, όχι όμως για να θεωρηθεί ότι αναπαράχθηκε πλήρως η κλασική απόδοση του Korf σε μεγάλα τυχαία σώματα.

Οι μηχανικές μετρήσεις της μηχανής H48 (Παράρτημα~\ref{app:h48eng}) δείχνουν ότι μεγάλο μέρος του κόστους των επαναλαμβανόμενων ρηχών κλήσεων ήταν μη αλγοριθμικό κόστος επικύρωσης και εκκίνησης: με δίκαιη λογιστική κοινής προθέρμανσης, η χρήση έμπιστου πίνακα, η λειτουργία δέσμης και η παραμένουσα διεργασία αφαιρούν αυτό το κόστος χωρίς να αλλάζουν την ίδια την αναζήτηση. Αντίθετα, οι δύσκολες μεμονωμένες περιπτώσεις εξακολουθούν να δαπανούν τον περισσότερο χρόνο στην ίδια την αναζήτηση H48.

Για το \eng{superflip}, η απόσταση 20 αποδεικνύεται πλήρως με εγγενή κώδικα. Η επαληθευμένη λύση μήκους 20 παρέχει το άνω φράγμα, ενώ ο εξαντλητικός IDA* ενιαίου ορίου τύπου Reid, με συμμετρικά ανηγμένο πίνακα φάσης~1 και παραδεκτό κάτω φράγμα τριών αξόνων \cite{reid1997optimal}, αποκλείει κάθε λύση μήκους 19 ή μικρότερου. Η απόδειξη αυτή ολοκληρώνεται σε περίπου 87 λεπτά σε συμβατικό μηχάνημα 8 πυρήνων/16~GiB και παρουσιάζεται αναλυτικά στο Κεφάλαιο~\ref{chap:experiments}. Το αποτέλεσμα συμφωνεί με το βιβλιογραφικά γνωστό κάτω φράγμα του Reid \cite{reid1995superflip}, ενώ η πιστοποίηση μέσω H48/Nissy λειτουργεί ως ανεξάρτητη διασταύρωση και όχι ως η μόνη πηγή.

Η πλήρης μελέτη του κύβου 2×2×2 λειτουργεί ως έλεγχος ωριμότητας της πειραματικής μεθοδολογίας: πλήρης απαρίθμηση, σαφώς ορισμένος χώρος καταστάσεων, αποθηκευμένη κατανομή αποστάσεων και αντιπροσωπευτικές ακριβείς λύσεις. Δεν αποτελεί υπεκφυγή από τον στόχο του 3×3×3, αλλά υποστηρικτικό παράδειγμα πλήρους απαρίθμησης σε χώρο που χωρά σε συμβατικό υπολογιστή. Η κατανομή της αφορά το συγκεκριμένο μοντέλο σταθερής αναφοράς, με σταθερή τη γωνία DBL και σύνολο κινήσεων \(U,R,F\), και όχι κάθε δυνατό φυσικό προσανατολισμό με όλες τις συμμετρίες. Η ίδια αρχή διατρέχει ολόκληρη την εργασία: κάθε αποτέλεσμα είναι ισχυρό μόνο μέσα στο μοντέλο στο οποίο δηλώθηκε.

Η κριτική αποτίμηση ολοκληρώνεται με τις απειλές εγκυρότητας των ευρημάτων, οι οποίες είναι τέσσερις. Πρώτη είναι η μικρή κλίμακα του σώματος μετρήσεων: οι 13 περιπτώσεις ανά επιλυτή αρκούν για αναπαραγώγιμη επίδειξη, όχι για στατιστική γενίκευση της απόδοσης. Δεύτερη είναι η τοπική φύση των μετρήσεων χρόνου, οι οποίες επηρεάζονται από το μηχάνημα, το περιβάλλον εκτέλεσης της Python και τα φαινόμενα κρυφής μνήμης. Τρίτη είναι η περιορισμένη εμβέλεια των εγγενών επιλυτών, οι οποίοι λύνουν το σώμα μετρήσεων χωρίς να αποτελούν καθολική βέλτιστη μέθοδο για τον κύβο 3×3×3. Τέταρτη είναι η προέλευση της μηχανής H48, η οποία βασίζεται σε δημόσιο επιλυτή και δεν αποτελεί πρωτότυπο αλγόριθμο της εργασίας, όπως ορίζεται στην Ενότητα~\ref{sec:intro-contributions}.

Η εργασία μετριάζει τις απειλές αυτές με διαφάνεια. Για την πρώτη, αποθηκεύεται κάθε γραμμή αποτελέσματος, κάθε ανακάτεμα, κάθε ετικέτα κατάστασης και κάθε αποτέλεσμα επαλήθευσης. Για τη δεύτερη, οι χρόνοι παρουσιάζονται ως ενδεικτικές μετρήσεις και όχι ως καθολικοί ισχυρισμοί απόδοσης. Για την τρίτη, χρησιμοποιούνται οι ετικέτες κατάστασης και οι ρητά διατυπωμένοι περιορισμοί της Ενότητας~\ref{sec:conclusions-limitations}. Για την τέταρτη, καταγράφονται άδεια, παραπομπές, αθροίσματα ελέγχου και μεταδεδομένα προέλευσης, όπως περιγράφεται στο Κεφάλαιο~\ref{chap:implementation}.

\section{Απάντηση στους στόχους της εργασίας}
\label{sec:conclusions-goals}

Το επίσημο κείμενο ανάθεσης ζητά να επιτευχθούν όσο το δυνατόν περισσότεροι από τους τρεις στόχους του: η υλοποίηση των αλγορίθμων Thistlethwaite, Kociemba και Korf, η αναγνώριση της απόστασης μιας κατάστασης από τη λυμένη θέση και η εύρεση κατάλληλης ευρετικής συνάρτησης για βέλτιστη επίλυση με τον αλγόριθμο A* ή παραλλαγή του (Κεφάλαιο~\ref{chap:intro}). Και οι τρεις στόχοι υλοποιούνται και τεκμηριώνονται με αποθηκευμένα τεκμήρια, με ρητά δηλωμένο το εύρος κάθε ισχυρισμού.

Ως προς τον πρώτο στόχο, η εργασία υλοποιεί περιορισμένης εμβέλειας εκδοχές και των τριών αλγορίθμων (Κεφάλαια~\ref{chap:algorithms} και~\ref{chap:implementation}). Ο τετραφασικός επιλυτής Thistlethwaite, στηριγμένος σε τέσσερις ακριβείς φασικούς πίνακες αποστάσεων, τερματίζει εγγυημένα σε κάθε έγκυρη κατάσταση. Ο επιλυτής Kociemba λύνει με τις προεπιλογές του τυπικά τυχαία ανακατέματα 18--20 κινήσεων. Ο επιλυτής Korf/IDA* αποδεικνύει ακριβείς αποστάσεις όπου η αναζήτηση ολοκληρώνεται. Οι κρίσιμοι για την απόδοση μηχανισμοί υλοποιούνται σε εγγενή κώδικα C/C++ που ενορχηστρώνεται από την Python, ώστε η εκτέλεση να παραμένει εφικτή σε συμβατικό υπολογιστή με 16~GiB μνήμης, όπως απαιτεί το κείμενο ανάθεσης.

Ως προς τον δεύτερο στόχο, η υλοποίηση δέχεται κατάσταση κύβου και επιστρέφει μία από τις κατηγορίες \code{exact_distance}, \code{lower_bound}, \code{unknown_timeout} ή \code{invalid_state}, ώστε η πρακτική επίλυση να μη συγχέεται με αποδεδειγμένη απόσταση. Η αναγνώριση απόστασης απαντάται δηλαδή με κατηγορίες τεκμηρίωσης και όχι με ανεπιφύλακτο καθολικό ισχυρισμό χρόνου εκτέλεσης· τα αντίστοιχα πειράματα παρουσιάζονται στο Κεφάλαιο~\ref{chap:experiments}.

Ως προς τον τρίτο στόχο, η απάντηση είναι ο επιλυτής Korf/IDA*, δηλαδή η επιλεγμένη παραλλαγή του A*, με παραδεκτά κάτω φράγματα από πίνακες: την πλήρη γωνιακή PDB, τους οκτώ εξακμικούς πίνακες προτύπων και τα χωριστά πειράματα επιμερισμένου κόστους. Τα τεκμήρια της μηχανής H48 (οι 15/15 ακριβείς γραμμές καταπόνησης άμεσης κατάστασης του \code{h48h7} και η πιστοποίηση του \eng{superflip}) λειτουργούν συμπληρωματικά και δεν υποκαθιστούν την εγγενή ευρετική Korf/IDA*. Το όριο των ισχυρισμών ελέγχεται και μηχανικά από το τεκμήριο συμβολαίου του Παραρτήματος~\ref{app:h48eng}: η διεπαφή του βέλτιστου μαντείου (\eng{oracle}) είναι υλοποιημένη για κάθε φυσικά έγκυρη κατάσταση 3×3×3 και το εμπειρικό σώμα γρήγορων περιπτώσεων υποστηρίζεται· ο ίδιος έλεγχος, ωστόσο, καταγράφει ότι δεν υπάρχει απόδειξη γρήγορου χρόνου εκτέλεσης για κάθε κατάσταση.

Η φράση «όσο το δυνατόν περισσότεροι» του κειμένου ανάθεσης δεν επιτρέπει χαλαρή κάλυψη χωρίς αλγοριθμική ουσία, αλλά ούτε και την παρουσίαση μερικών επιλυτών ως πλήρων. Η εργασία δεν περιορίζεται σε περιβλήματα έτοιμων εργαλείων: περιλαμβάνει εγγενή φασικά κατηγορήματα, πίνακες συντεταγμένων, πίνακες αποκοπής (\eng{pruning tables}), εγγενή γωνιακή PDB του 3×3×3 και την εξαντλητική μελέτη του 2×2×2. Η ισορροπία της βρίσκεται ακριβώς σε αυτή τη γραμμή.

\section{Περιορισμοί}
\label{sec:conclusions-limitations}

Οι περιορισμοί που ακολουθούν αποτελούν μέρος της επιστημονικής συνεισφοράς και δηλώνονται με την ίδια αυστηρότητα με τα θετικά αποτελέσματα.

Πρώτον, η κάλυψη της κλασικής βιβλιογραφίας βέλτιστης επίλυσης είναι περιορισμένη. Η υποδομή με πλήρη γωνιακή PDB, οκτώ εξακμικούς πίνακες προτύπων και δύο πειραματικούς πίνακες ακμών επιμερισμένου κόστους συνιστά ουσιαστικό βήμα πέρα από μια καθαρά θεωρητική περιγραφή ή μια απλή ευρετική σε Python. Δεν ισοδυναμεί όμως με πλήρεις πολλαπλούς προσθετικούς πίνακες προτύπων, ούτε με αναπαραγωγή μεγάλης κλίμακας των ιστορικών πειραμάτων Korf ή H48.

Δεύτερον, στον επιλυτή Korf/IDA* καθαρής Python παραμένουν εξαντλήσεις χρονικού ορίου για τις βαθύτερες τυχαίες περιπτώσεις του σώματος μετρήσεων, ενώ οι δύσκολες μεμονωμένες αναζητήσεις H48 παραμένουν χρονικά δαπανηρές ακόμη και με έμπιστο πίνακα και προφόρτωση. Η ζωντανή πιστοποίηση βελτιστότητας μέσα στη δοκιμαστική σουίτα φτάνει έως την απόσταση 6 για την υλοποίηση καθαρής Python και έως την απόσταση 11 μέσω του αμοιβαίου ελέγχου, κατά τον οποίο ο εγγενής επιλυτής Korf και η μηχανή απόδειξης βελτιστότητας δύο φάσεων αποδεικνύουν ανεξάρτητα το ίδιο μήκος. Πέρα από αυτά τα βάθη, η ορθότητα στηρίζεται στις ολοκληρωμένες παραδεκτές αναζητήσεις και στις διασταυρώσεις του Κεφαλαίου~\ref{chap:validation}.

Τρίτον, ο πίνακας \code{h48h7} δεν αναπαράγεται πανομοιότυπος byte προς byte, επειδή η πολυνηματική παραγωγή της περιόδου δημιουργίας του ήταν μη ντετερμινιστική· τα τεκμήρια πιστοποίησης που τον συνοδεύουν διατηρούνται ως τεκμηριωμένη απόφαση διατήρησης. Η εγκυρότητα των βαθιών ακριβών γραμμών δεν επαφίεται πάντως σε αυτόν τον πίνακα, αφού όλες οι γραμμές που στηρίζονταν αποκλειστικά σε αυτόν έχουν επιβεβαιωθεί ανεξάρτητα από το εξωτερικό Nissy 2.0.8, όπως τεκμηριώνεται στο Κεφάλαιο~\ref{chap:validation}.

Τέταρτον, η εργασία δεν διατυπώνει τυπικό θεώρημα πρακτικού χρόνου εκτέλεσης για ολόκληρο τον χώρο καταστάσεων του 3×3×3. Οι σχετικοί ισχυρισμοί παραμένουν εμπειρικοί και περιορισμένοι στο μετρημένο σώμα, όπως αποτυπώνεται και στο τεκμήριο συμβολαίου που αναφέρθηκε στην Ενότητα~\ref{sec:conclusions-goals}.

\section{Μελλοντικές Βελτιώσεις}
\label{sec:conclusions-future}

Οι κατευθύνσεις βελτίωσης που προκύπτουν από τα αποτελέσματα συνοψίζονται σε έναν ενιαίο κατάλογο:

\begin{itemize}
\item \textbf{Ισχυρότερες συντεταγμένες δεύτερης φάσης για τον επιλυτή Kociemba}: μια συνδυασμένη φασική συντεταγμένη που ενσωματώνει περισσότερη πληροφορία μεταθέσεων θα έκανε την εγγενή υλοποίηση σταθερότερη σε μεγαλύτερα σώματα και σε χαμηλότερα όρια χρόνου.
\item \textbf{Πίνακες ελιγμών και βελτιστοποίηση φασικών ορίων για τον επιλυτή Thistlethwaite}: αφού η πληρότητα είναι ήδη εγγυημένη, το επόμενο βήμα αφορά το μήκος των λύσεων, με ιστορικούς πίνακες ελιγμών (\eng{maneuver databases}) και αναζήτηση που εξετάζει εναλλακτικές ισομήκεις φασικές λύσεις ή ανέχεται προσωρινά μη φθίνουσα φασική απόσταση όταν μειώνει το συνολικό μήκος.
\item \textbf{Γενίκευση του προσθετικού ή επιμερισμένου κόστους συνδυασμού πινάκων προτύπων} πέρα από τα περιορισμένα πειράματα URF/DLB, με αποδεδειγμένη παραδεκτότητα, αναγωγή συμμετρίας και καλύτερη αξιοποίηση των υπαρχόντων εξακμικών πινάκων στον επιλυτή Korf/IDA*.
\item \textbf{Μεγαλύτεροι πίνακες H48} (\code{h48h8}, \code{h48h9}), εφόσον γίνει αποδεκτό το κόστος παραγωγής και αποθήκευσης· η τοπική δοκιμή παραγωγής του \code{h48h8} έδειξε ότι το επόμενο επίπεδο είναι κυρίως ζήτημα εγγενούς παραγωγής πινάκων, μνήμης και εισόδου/εξόδου για συμβατικό μηχάνημα 16~GiB, όχι βελτιστοποίησης της Python.
\item \textbf{Επιτάχυνση της εγγενούς εξαντλητικής απόδειξης του \eng{superflip}}, για παράδειγμα με πίνακες συμπιεσμένους σε στοιχεία των 4 bit (\eng{nibbles}) και σταδιακή αποκοπή, ώστε η ίδια απόδειξη να ολοκληρώνεται σε λίγα λεπτά αντί για περίπου 87 λεπτά· το περιθώριο βελτίωσης αφορά τον χρόνο εκτέλεσης και όχι την ορθότητα.
\item \textbf{Μεγαλύτερο σώμα μετρήσεων και καταπόνησης}, με επαναλήψεις των μετρήσεων χρόνου και τυπική τεκμηρίωση του ορίου ισχυρισμών για όλες τις καταστάσεις.
\end{itemize}

Κάθε τέτοια επέκταση πρέπει να συνοδεύεται από δοκιμές, μεταδεδομένα παραγωγής και νέα αποθηκευμένα τεκμήρια, ώστε οι ισχυρισμοί του κειμένου να ενισχύονται μόνο όσο επιτρέπει η νέα απόδειξη. Αντίστροφα, όπου δεν προστίθεται ισχυρότερη απόδειξη, οι περιορισμοί πρέπει να παραμένουν ορατοί.

\section{Επίλογος}
\label{sec:conclusions-epilogue}

Η βασική συνεισφορά της παρούσας εργασίας δεν είναι ο ισχυρισμός ότι κάθε αυθαίρετη κατάσταση 3×3×3 λύνεται αποδεδειγμένα βέλτιστα σε μικρό χρόνο χειρότερης περίπτωσης· κάτι τέτοιο ούτε αποδεικνύεται ούτε δηλώνεται. Η συνεισφορά είναι ένα συνεκτικό, ελέγξιμο πλαίσιο στο οποίο οι διαφορετικές μορφές επίλυσης διαχωρίζονται καθαρά: ακριβείς απαντήσεις όπου η παραδεκτή αναζήτηση ολοκληρώνεται, επαληθευμένες πρακτικές λύσεις όπου οι σταδιακές μέθοδοι επιτυγχάνουν, και ρητά κάτω φράγματα ή εξαντλήσεις χρονικού ορίου παντού αλλού. Το πλαίσιο συμπληρώνεται από μια πρακτική μεθοδολογία συγχρονισμού κώδικα και κειμένου: οι αριθμοί των πινάκων δεν πληκτρολογούνται με το χέρι, αλλά παράγονται από σενάρια, αποθηκεύονται σε αρχεία αποτελεσμάτων, επαληθεύονται και εισάγονται στο κείμενο (Κεφάλαιο~\ref{chap:validation}, Παράρτημα~\ref{app:repro}). Η τελική αξία της εργασίας δεν εξαρτάται μόνο από το πόσες περιπτώσεις λύνονται, αλλά και από το ότι ο αναγνώστης μπορεί να εμπιστευθεί ότι κάθε ετικέτα, κάθε πίνακας και κάθε συμπέρασμα αντιστοιχούν σε πραγματικό, αποθηκευμένο τεκμήριο.
